本附录收录正文中引用但因篇幅置于此处的完整表格。全部表格由 \nolinkurl{tables} 阶段生成
（\nolinkurl{.tex} 与同名 \nolinkurl{.csv} 并存），数值与结果工作簿一致。

\subsection{问题 1：冲突强度、度分布与全部冲突对}

\begin{table}[H]
\centering\small
\begin{minipage}[t]{0.46\linewidth}
\centering
\caption{冲突强度分布：横跨 $k$ 对使用次的冲突对数量。多数冲突对只在一次使用上相撞。}\label{tab:q1mult}
\input{generated/tables/tab_q1_multiplicity.tex}
\end{minipage}\hfill
\begin{minipage}[t]{0.46\linewidth}
\centering
\caption{冲突图中度最大的 $10$ 个用频装备。}\label{tab:q1deg}
\input{generated/tables/tab_q1_top_degree.tex}
\end{minipage}
\end{table}


\needspace{12\baselineskip}
\captionof{table}{问题 1 检测出的全部 $\valQoneConflictPairs$ 对时频冲突（按装备编号字典序，与 \nolinkurl{result1.xlsx} 一致）。
“冲突使用次对数”为该计划对中发生时间重叠的使用次对数，“频段重叠宽度”以 $\Delta f$ 计。}
\label{tab:q1conflicts}
{\footnotesize\input{generated/tables/tab_q1_conflicts.tex}}
\addtocounter{table}{-1}% longtable 自身也步进了 table 计数器，此处抵消，保证表号连续

\subsection{问题 2：被撤销计划与逐计划动作}

\begin{table}[H]
\centering\small
\caption{问题 2 中被撤销计划的局部环境。“冲突邻居数”为该计划在冲突图中的度，
“其中 A/B 类”为这些邻居中高优先级装备的个数，“可用平移动作数”为式~\eqref{eq:options} 中该计划的合法平移动作数。
多数被撤销的计划被高优先级邻居包围；表中也有邻居数仅为 $1$ 的计划，说明原始度数不是撤销的直接原因，
真正起作用的是最后一列——没有任何一个合法单参数动作能使它与存留计划都不冲突。}
\label{tab:q2cancel}
\resizebox{\linewidth}{!}{\input{generated/tables/tab_q2_cancelled.tex}}
\end{table}

\needspace{12\baselineskip}
\captionof{table}{问题 2 的逐计划动作清单（只列出被调整或被撤销的计划，与 \nolinkurl{result2.xlsx} 一致）。}
\label{tab:q2adj}
{\footnotesize\input{generated/tables/tab_q2_adjustments.tex}}
\addtocounter{table}{-1}% longtable 自身也步进了 table 计数器，此处抵消，保证表号连续

\subsection{问题 3：新增计划清单}

\needspace{12\baselineskip}
\captionof{table}{问题 3 在解释 A 下新增的 $\valQthreeNewPlans$ 个 C 类用频计划（与 \nolinkurl{result3.xlsx} 一致；
间隔时长与使用次数按假设~\ref{as:pack} 与附件 C 类相同）。}
\label{tab:q3plans}
{\footnotesize\input{generated/tables/tab_q3_plans.tex}}
\addtocounter{table}{-1}% longtable 自身也步进了 table 计数器，此处抵消，保证表号连续

\subsection{问题 4：模型规模、目标方案与逐计划动作}

\begin{table}[H]
\centering\small
\begin{minipage}[t]{0.52\linewidth}
\centering
\caption{问题 4 的模型规模与求解设置。}\label{tab:q4model}
\resizebox{\linewidth}{!}{\input{generated/tables/tab_q4_model.tex}}
\end{minipage}\hfill
\begin{minipage}[t]{0.44\linewidth}
\centering
\caption{问题 4 中被撤销计划的局部环境（列的含义同表~\ref{tab:q2cancel}）。}\label{tab:q4cancel}
\resizebox{\linewidth}{!}{\input{generated/tables/tab_q4_cancelled.tex}}
\end{minipage}
\end{table}

\begin{table}[H]
\centering\small
\caption{四种目标方案在问题 4 上的对照（列的含义同表~\ref{tab:q2schemes}，另给出使用间隔调整动作的计划数）。}
\label{tab:q4schemes}
\resizebox{\linewidth}{!}{\input{generated/tables/tab_q4_schemes.tex}}
\end{table}

\needspace{12\baselineskip}
\captionof{table}{问题 4 的逐计划动作清单（只列出被调整或被撤销的计划，与 \nolinkurl{result4.xlsx} 一致）。}
\label{tab:q4adj}
{\footnotesize\input{generated/tables/tab_q4_adjustments.tex}}
\addtocounter{table}{-1}% longtable 自身也步进了 table 计数器，此处抵消，保证表号连续

\subsection{模型检验：微型实例穷举对照与规模基准}

\begin{table}[H]
\centering\footnotesize
\caption{$\valBenchTinyInstances$ 个随机微型实例上的穷举对照。对每个实例枚举全部动作组合，
求出词典序最优向量，与本文的 CP-SAT 求解器结果逐一比对；最大枚举组合数为 $\valBenchTinyMaxCombinations$。}
\label{tab:benchtiny}
\setlength{\tabcolsep}{3pt}
{\scriptsize\input{generated/tables/tab_bench_tiny.tex}}
\end{table}

\begin{table}[H]
\centering\footnotesize
\caption{合成实例上的规模基准（$n=\valBenchScaleMinN,\dots,\valBenchScaleMaxN$，时间范围随 $n$ 等比放大以保持密度）。
CP-SAT 一列采用固定 $120$ s 时限，故其状态只说明“在该预算内找到了可行解”，不代表可证明最优。}
\label{tab:benchscaling}
\setlength{\tabcolsep}{4pt}
\resizebox{\linewidth}{!}{\input{generated/tables/tab_bench_scaling.tex}}
\end{table}

\begin{figure}[H]
\centering
\includegraphics[width=0.92\linewidth]{fig_bench_scaling}
\caption{算法的规模基准。合成实例 $n=\valBenchScaleMinN,\dots,\valBenchScaleMaxN$ 下
两种检测算法与模型构建的用时，以及 CP-SAT 在固定 $120$ s 时限下的求解时间（纵轴为对数刻度）。
检测算法始终为亚秒级，求解时间则受时限截断。}
\label{fig:benchscaling}
\end{figure}

\FloatBarrier

\subsection{最优性证据的两张补充表}

\begin{table}[H]
\centering\footnotesize
\caption{问题 2 的条件最优性（MDR-0012）：把被撤销计划的\emph{集合}固定为表~\ref{tab:q2table1} 所报告的那
$\valQtwoValueCancelc$ 个 C 类计划之后，对调整三级与幅度级逐级最小化的结果。此处的“已证明下界”与状态
\textsf{OPTIMAL} 只对该\emph{受限}问题成立，不构成对词典序真优解相应分量的界。}
\label{tab:q2conditional}
\input{generated/tables/tab_q2_conditional.tex}
\end{table}

\begin{table}[H]
\centering\small
\caption{问题 3 的答案对“采用哪一个问题 2 方案”的依赖。四个方案均为无冲突方案且均通过独立校验；
在每个方案的存留计划布局之上求解同一个模型~\eqref{eq:pack}。各方案的时间范围可能不同，故单列给出，
不同布局的自由单元数不可直接相减。}\label{tab:q3layouts}
\resizebox{\linewidth}{!}{\input{generated/tables/tab_q3_layouts.tex}}
\end{table}
